\section*{Chapter \ref{ch:b3:rovibrational-spectroscopy} --- Rotational and Vibrational Spectroscopy}

\begin{solution}{exo:b3:rovibrational-spectroscopy:1}
$I = h/(8\pi^2cB) = 6.62607 \times 10^{-34}/(8\pi^2 \times 2.99792 \times 10^{10} \times
1.9225) = \qty{1.4561e-46}{kg.m^2}$. $\mu = \qty{6.85621}{u} = \qty{1.13850e-26}{kg}$, so
$r_0 = \sqrt{I/\mu} = \qty{113.09}{pm}$.
\end{solution}

\begin{solution}{exo:b3:rovibrational-spectroscopy:2}
$2B_0(J+1)$: 20.88, 41.76, 62.64, \qty{83.52}{cm^{-1}}. $J_{\max} = \sqrt{kT/2hcB} -
\frac12 = \sqrt{208.5/20.88} - 0.5 = 2.7$: the level $J = 3$ is the most
populated, and the strongest line is $4 \leftarrow 3$ at \qty{83.5}{cm^{-1}}.
\end{solution}

\begin{solution}{exo:b3:rovibrational-spectroscopy:3}
Microwave absorption needs a permanent dipole: \ce{HCl} and \ce{H2O} only.
Rotational Raman needs an anisotropic polarisability: \ce{H2}, \ce{HCl}, \ce{N2},
\ce{CO2}, \ce{H2O}; not the \omterm{def:b3:rovibrational-spectroscopy:rotor-types}{spherical tops} \ce{CH4} and \ce{SF6}.
\end{solution}

\begin{solution}{exo:b3:rovibrational-spectroscopy:4}
$N_1/N_0 = \eu^{-hc\tilde\nu/kT}$: at \qty{300}{K}, $\eu^{-13.84} = \num{9.8e-7}$; at
\qty{1000}{K}, $\eu^{-4.15} = 0.016$. The \omterm{def:b3:rovibrational-spectroscopy:anharmonic}{hot band} becomes noticeable (a per cent of
the fundamental) only near \qty{1000}{K}.
\end{solution}

\begin{solution}{exo:b3:rovibrational-spectroscopy:5}
$\tilde\omega_e - 2\tilde\omega_ex_e = 2885.3$ and $2\tilde\omega_e - 6\tilde\omega_ex_e =
5665.0$: subtracting twice the first from the second, $-2\tilde\omega_ex_e = -105.6$,
so $\tilde\omega_ex_e = \qty{52.8}{cm^{-1}}$ and $\tilde\omega_e = \qty{2990.9}{cm^{-1}}$.
\end{solution}

\begin{solution}{exo:b3:rovibrational-spectroscopy:6}
$D_e = 2990.95^2/(4 \times 52.82) = \qty{42341}{cm^{-1}}$; $G(0) = 1495.47 - 13.20 =
\qty{1482.3}{cm^{-1}}$; $D_0 = \qty{40859}{cm^{-1}} = \qty{488.8}{kJ/mol}$, against the true
\qty{35760}{cm^{-1}} $= \qty{427.8}{kJ/mol}$: the Morse model overestimates by 14\,\%.
\end{solution}

\begin{solution}{exo:b3:rovibrational-spectroscopy:7}
$R(0) - P(2) = 6B_0 = 62.64$: $B_0 = \qty{10.440}{cm^{-1}}$. $R(1) - P(1) = 6B_1 = 60.80$:
$B_1 = \qty{10.133}{cm^{-1}}$. $\alpha_e = \qty{0.307}{cm^{-1}}$. $R(0) = \tilde\nu_0 + 2B_1$ gives
$\tilde\nu_0 = \qty{2885.31}{cm^{-1}}$ (and $P(1) = \tilde\nu_0 - 2B_0$ agrees).
\end{solution}

\begin{solution}{exo:b3:rovibrational-spectroscopy:8}
$\mu_H = \qty{0.97959}{u}$, $\mu_D = \qty{1.90441}{u}$, $\mu_H/\mu_D = 0.51438$. $B_0(\ce{DCl})
= 10.440 \times 0.51438 = \qty{5.370}{cm^{-1}}$. $\tilde\omega_e = 2990.95 \times \sqrt{0.51438}
= 2145.1$, $\tilde\omega_ex_e = 52.82 \times 0.51438 = 27.17$: $\tilde\nu_0 = 2145.1 - 54.3 =
\qty{2090.8}{cm^{-1}}$.
\end{solution}

\begin{solution}{exo:b3:rovibrational-spectroscopy:9}
Shifts $4B_0(J + \frac32)$: 11.94, 19.90, \qty{27.85}{cm^{-1}} (from $J = 0$, 1, 2). The
two \ce{^{14}N} nuclei are identical bosons of spin 1: the even-$J$ levels have six
nuclear-spin states, the odd-$J$ three, so lines from even $J$ are twice as
strong as their neighbours.
\end{solution}

\begin{solution}{exo:b3:rovibrational-spectroscopy:10}
A \omterm{def:b3:quantum-model-systems:rotor}{rigid rotor} would put the lines at $\nu_1$, $2\nu_1$, $3\nu_1$: $230542.40$ and
\qty{345813.61}{MHz}, above the measured ones by 4.4 and \qty{17.6}{MHz}. With
$\nu_{J+1\leftarrow J} = 2B(J+1) - 4D(J+1)^3$: $\nu_1 = 2B - 4D$, $\nu_3 = 6B - 108D$, so
$3\nu_1 - \nu_3 = 96D$: $D = \qty{0.1835}{MHz}$ and $B = (\nu_1 + 4D)/2 =
\qty{57635.97}{MHz}$. Check: $\nu_2 = 4B - 32D = \qty{230538.0}{MHz}$.
\end{solution}

\begin{solution}{exo:b3:rovibrational-spectroscopy:11}
$\Delta G(v + \frac12) = \tilde\omega_e - 2\tilde\omega_ex_e(v+1)$ vanishes near $v =
\tilde\omega_e/2\tilde\omega_ex_e - 1$; the sum of an arithmetic progression of $n$ terms
from $\tilde\omega_e - 2\tilde\omega_ex_e$ down to almost 0 is about $n\tilde\omega_e/2 \approx
\tilde\omega_e^2/4\tilde\omega_ex_e - \tilde\omega_e/2 \approx D_e - G(0)$. For \ce{HCl}, the 28 gaps
($v = 0$ to 27) add up to \qty{40857}{cm^{-1}}, against $D_e - G(0) = \qty{40859}{cm^{-1}}$.
\end{solution}

\begin{solution}{exo:b3:rovibrational-spectroscopy:12}
With only even $J$, the transitions $J \to J + 2$ start at $J = 0, 2, 4, \dots$: shifts
$6B, 14B, 22B, \dots$, spacing $8B$. \ce{CO2} has a centre of symmetry and no
permanent dipole: no microwave spectrum.
\end{solution}

\begin{solution}{pb:b3:rovibrational-spectroscopy:1}
\textbf{1.} $B_0 = \nu/2 = \qty{57.6356}{GHz}$; $B_0/c = \qty{1.92252}{cm^{-1}}$.
\textbf{2.} $\mu = 12 \times 15.99491/27.99491 = \qty{6.85621}{u} = \qty{1.13850e-26}{kg}$.
\textbf{3.} $I = h/8\pi^2B_0 = \qty{1.4560e-46}{kg.m^2}$.
\textbf{4.} $r_0 = \sqrt{I/\mu} = \qty{113.09}{pm}$.
\textbf{5.} $r_0$ exceeds $r_e = \qty{112.83}{pm}$ by \qty{0.26}{pm}: $B_0$ averages $1/r^2$
over the zero-point vibration in an anharmonic well, which spends more time
at long distances.
\textbf{6.} $2\nu_{1\leftarrow0} = \qty{230.5424}{GHz}$; the measured line is \qty{4.4}{MHz} lower:
centrifugal distortion.
\textbf{7.} $0$, 3.845 and \qty{11.535}{cm^{-1}}, i.e. $E/k = 0$, 5.53 and \qty{16.60}{K}.
\textbf{8.} $N_2/N_1 = \frac53\exp[-(16.60 - 5.53)\,\mathrm K/T]$.
\textbf{9.} $\ln(\frac53/0.85) = 0.673 = 11.06/T$: $T = \qty{16}{K}$.
\textbf{10.} $J_{\max} = \sqrt{16.4/(2 \times 1.43878 \times 1.9225)} - 0.5 = 1.2$: $J = 1$.
\textbf{11.} A vibrational quantum corresponds to $hc\tilde\nu/k \approx \qty{3080}{K}$: at
\qty{16}{K} no molecule is vibrationally excited, so the cloud emits no vibrational
light.
\textbf{12.} It needs only \qty{5.5}{K} of excitation, is excited even in the coldest
gas, and CO is abundant and has a dipole moment.
\textbf{13.} $\mu = 13.00335 \times 15.99491/28.99826 = \qty{7.17241}{u}$.
\textbf{14.} $B \propto 1/\mu$: $115.2712 \times 6.85621/7.17241 = \qty{110.1893}{GHz}$.
\textbf{15.} The prediction is \qty{12}{MHz} (0.011\,\%) below the measured line: $B_0 = B_e - \frac12\alpha_e$, and the
vibration--rotation term $\alpha_e$ scales with a different power of $\mu$; the
equilibrium structure is the same, the zero-point average is not.
\textbf{16.} $98.9/1.1 = 90$.
\textbf{17.} The \ce{^{12}C^{16}O} line is saturated (optically thick): it no longer
grows with the amount of gas.
\textbf{18.} The rare \ce{^{13}C^{16}O} line stays proportional to the amount of gas, the
\ce{^{12}C^{16}O} line gives the temperature: together they measure both.
\textbf{19.} $\tilde\nu_0 = 2169.81 - 2 \times 13.288 = \qty{2143.24}{cm^{-1}}$, \qty{4.666}{\micro m}.
\textbf{20.} $2\tilde\omega_e - 6\tilde\omega_ex_e = \qty{4259.89}{cm^{-1}}$.
\textbf{21.} $G(0) = \frac12\tilde\omega_e - \frac14\tilde\omega_ex_e = \qty{1081.58}{cm^{-1}}$.
\textbf{22.} $R(0) - P(1) = 2(B_0 + B_1) \approx 4B \approx \qty{7.7}{cm^{-1}}$.
\textbf{23.} \ce{^{13}C^{16}O}: heavier \omterm{def:b3:quantum-model-systems:oscillator}{reduced mass}, lower frequency.
\textbf{24.} \textbf{$r_0(\ce{CO}) = \qty{113.1}{pm}$}, from a single radio frequency.
\end{solution}
